\section{MATLAB data extraction scripts}
Thanks to the ton of already available libraries and functions for data processing, MATLAB has been the first choice when dealing with the elaboration of all the sensor information retrieved from the data collection campaigns. Transforming CSV values into MATLAB data structure is immediate (the \textit{readtable} function is enough), and then the Signal Analysis Toolbox contains an extensive set of tools to work on those values with ease. In particular, I prepared the following scripts:
\begin{itemize} 
    \item \textbf{Data Extraction for ST devices}: works with Acceleration, Gyroscope and Magnetometer data. It reads all the values, and plots them over time using the MATLAB plotting feature. Then, all the data is modified for the Signal Analyzer application: the timesync set of values is transformed into a square wave, to represent the moments in which the re-synchronization happens, then all the signals are coupled with their timestamp generating \textit{timeseries} data structures. Those timeseries are further analyzed, to order them in the time domain and remove duplicates. In this way, the signals are ready to be shown and elaborated with MATLAB toolboxes.
    \item \textbf{Data Extraction for Meta devices}: it behaves like the one for ST devices, but takes into account the slight differences of CSV data generated by the MBIENTLAB ones. Then, everything remains the same: plots are shown, \textit{timeseries} are generated and elaborated. In particular, the timeseries variables have been called in a different way, in order for them to be distinguished from the ST ones.
    \item \textbf{Data Extract for Voice}: the goal of this script is to reconstruct the audio signal from the CSV file generated by the BlueTile with BlueVoice enabled. As the timestamping mechanism is not as precise as the one used for sensor data (due to the closed-source BlueVoice library), only the timestamps generated at resynchronization events are used, and all the samples in the middle receive regularly intervalled timestamps between the two events. Reducing the re-synchronization interval could be advantageous in this case, but the number of packet loss could be higher, making the audio quality sensibly lower. Also in this case \textit{timeseries} structures are created and elaborated to be shown afterwards in the Signal Analyzed application.
    \item \textbf{Data Extract for Frequencies}: it works only with the frequencies in the CSV generated by the SensorTile.box device. It behaves like the acceleration/gyroscope script, but works with frequencies, that are transformed from the indexes between 0-128 into the real frequencies values, and then scaled by a factor of 1000: this is necessary to make them show in the same plot of the acceleration values, for example. Also in this case, everything is modified to be shown correctly in the Signal Analyzer application. 
\end{itemize}
